\section{Experimental Results}

\subsection{Dataset}

Experiments were conducted on a lung adenocarcinoma (LUAD) evolution cohort comprising matched single-nucleus RNA sequencing (snRNA-seq) and spatial transcriptomics data spanning five histopathological stages: Normal, atypical adenomatous hyperplasia (AAH), adenocarcinoma in situ (AIS), minimally invasive adenocarcinoma (MIA), and invasive LUAD \cite{Peng2025-xu}.

After quality control, the snRNA-seq dataset comprised 787,709 cells across 18,075 genes. The spatial transcriptomics data included matched Visium sections spanning all five histopathological stages.

Quality control filtering removed cells with fewer than 200 detected genes or greater than 20\% mitochondrial reads. Genes expressed in fewer than 10 cells were excluded. Expression values were normalized to counts per 10,000 and log-transformed. Batch effects across samples were corrected using Harmony integration prior to downstream analysis.


\subsection{Reference Model Preparation}

StageBridge requires dual-reference latent embeddings from both healthy (HLCA) and cancer (LuCA) atlas models. The LuCA scANVI model was retrained from scratch on the LuCA Core Atlas to ensure reproducibility and compatibility with the scArches surgical fine-tuning workflow.

\paragraph{LuCA scANVI Retraining.} The LuCA Core Atlas \cite{Salcher2022-luca} (892,296 cells, 21 datasets, 33 cell types) was used to train a scVI base model followed by scANVI label-aware fine-tuning. Highly variable gene selection identified 6,000 genes using the Seurat v3 method with batch-aware selection. The scVI encoder used 10 latent dimensions, 128 hidden units, and 2 layers with 0.2 dropout, trained for 329 epochs with early stopping (patience=20, monitoring validation ELBO). The scANVI model was initialized from scVI weights and trained for an additional 34 epochs with early stopping (patience=15).

\begin{table}[h]
\centering
\caption{LuCA scANVI model validation using scIB metrics \cite{Luecken2022-scib}. The retrained model matches or exceeds the original across all biological preservation metrics.}
\label{tab:luca_validation}
\begin{tabular}{lccc}
\toprule
\textbf{Metric} & \textbf{Retrained} & \textbf{Original} & \textbf{$\Delta$} \\
\midrule
Overall scIB Score & \textbf{0.673} & 0.669 & +0.6\% \\
Adjusted Rand Index (ARI) & \textbf{0.697} & 0.670 & +4.0\% \\
Normalized Mutual Information (NMI) & \textbf{0.821} & 0.816 & +0.6\% \\
Cell Type $k$-NN Purity & \textbf{0.922} & 0.918 & +0.4\% \\
Batch ASW & 0.555 & 0.575 & $-$3.5\% \\
Batch $k$-NN Mixing & \textbf{0.530} & 0.519 & +2.1\% \\
Cell Type ASW & \textbf{0.598} & 0.595 & +0.5\% \\
\bottomrule
\end{tabular}
\end{table}

The retrained model achieved comparable or improved performance across all biological preservation metrics (ARI, NMI, cell type purity) while maintaining batch integration quality. The slight decrease in batch ASW is offset by improved batch $k$-NN mixing, indicating the model learned to integrate batches through neighbor mixing rather than forcing global separation. Training was conducted on 4$\times$ NVIDIA H100 GPUs using PyTorch Lightning's \texttt{ddp\_spawn} distributed strategy, completing in approximately 6 hours. The model passed validation (overall scIB score 0.673 vs.\ original 0.669) and is ready for scArches surgery.


\subsection{Dual-Reference Mapping}

Query cells were projected onto both HLCA (healthy) and LuCA (cancer) reference latent spaces using scArches surgical fine-tuning. A key challenge was gene identifier reconciliation: the HLCA model expects 2,000 ENSG identifiers while the LuCA model expects 6,000 ENSG identifiers, but the query data uses gene symbols.

\paragraph{Gene Identifier Mapping.} To maximize gene coverage, we constructed a combined symbol-to-ENSG mapping from both references. The LuCA Core Atlas provided 17,764 symbol$\to$ENSG mappings via its \texttt{feature\_name} annotations, while the HLCA reference contributed 2,000 additional mappings. After deduplication (with HLCA taking precedence for conflicts), the combined mapping contained 17,865 unique symbols. This enabled conversion of 15,460/18,075 (85.5\%) query genes to ENSG identifiers.

\paragraph{HLCA Mapping.} The query achieved 86.05\% overlap with the HLCA model's 2,000 highly variable genes. scArches surgery was performed for up to 200 epochs with early stopping (patience=15, monitoring validation ELBO). Training converged at epoch 60--61, producing 30-dimensional latent embeddings for all 787,709 query cells.

\paragraph{LuCA Mapping.} Initial attempts with only HLCA-derived ENSG mappings yielded only 24.7\% overlap with the LuCA model's 6,000 genes, resulting in out-of-memory failures during gene padding. After incorporating LuCA's own symbol$\to$ENSG mapping, overlap improved substantially. scArches surgery was performed with identical hyperparameters, producing 10-dimensional latent embeddings.

\paragraph{Embedding Fusion.} Both reference embeddings were L2-normalized before concatenation to prevent scale domination, yielding 40-dimensional fused embeddings. Confidence scores were computed via percentile-rank calibration of $k$-NN distances to ensure comparability between references.

\begin{table}[h]
\centering
\caption{Dual-reference mapping summary. Gene overlap indicates percentage of model's HVGs found in query after ENSG conversion.}
\label{tab:dual_ref_mapping}
\begin{tabular}{lcccc}
\toprule
\textbf{Reference} & \textbf{Model Genes} & \textbf{Gene Overlap} & \textbf{Latent Dim} & \textbf{Epochs} \\
\midrule
HLCA & 2,000 & 86.05\% & 30 & 60--61 \\
LuCA & 6,000 & TBD & 10 & TBD \\
\midrule
\textbf{Fused} & -- & -- & 40 & -- \\
\bottomrule
\end{tabular}
\end{table}


\subsection{Spatial Mapping Backend Comparison}

Three spatial mapping backends were evaluated to determine the optimal method for projecting snRNA-seq cells onto tissue coordinates: Tangram \cite{Biancalani2021-bt}, DestVI \cite{Lopez2022-ps}, and TACCO. Each backend was assessed on mapping accuracy using held-out spatial spots with known cell type composition, cell type assignment quality via macro-averaged F1 score, spatial coherence measured by local Moran's I statistic, and computational cost measured as wall-clock runtime on standardized hardware.

\begin{table}[h]
\centering
\caption{Spatial backend comparison. Mapping accuracy measured on held-out spots; F1 is macro-averaged over cell types; Moran's I measures spatial autocorrelation of assigned types.}
\label{tab:spatial_backends}
\begin{tabular}{lcccc}
\toprule
\textbf{Backend} & \textbf{Mapping Acc.} & \textbf{Cell Type F1} & \textbf{Moran's I} & \textbf{Runtime (hrs)} \\
\midrule
Tangram & -- & -- & -- & -- \\
DestVI & -- & -- & -- & -- \\
TACCO & -- & -- & -- & -- \\
\bottomrule
\end{tabular}
\end{table}

\textbf{[FIGURE 1: Spatial backend comparison]}

\textit{Figure 1. Comparison of spatial mapping backends. (A) Cell type assignment accuracy on held-out spots. (B) Spatial coherence measured by local Moran's I. (C) Runtime scaling with dataset size. (D) Example spatial maps showing cell type assignments for each backend.}

\textbf{[RESULTS: Describe which backend performed best and why it was selected. Report accuracy, coherence, and runtime metrics. Justify backend choice for downstream analyses.]}


\subsection{Self-Supervised Pretraining}

The self-supervised pretraining stage was evaluated by monitoring convergence of the primary masked reconstruction objective and four auxiliary losses over training epochs. Validation metrics were tracked to assess generalization and detect potential overfitting. The learned representations were qualitatively examined via dimensionality reduction to assess whether biologically meaningful structure emerged prior to downstream fine-tuning.

\textbf{[FIGURE 2: SSL pretraining curves]}

\textit{Figure 2. Self-supervised pretraining convergence. (A) Total loss over training epochs. (B) Individual loss components (masked reconstruction, ranking, consistency, spatial, group). (C) Validation loss showing generalization. (D) Learning rate schedule with warmup and cosine decay.}

\begin{table}[h]
\centering
\caption{Self-supervised pretraining metrics. Initial values at epoch 1; final values at convergence.}
\label{tab:ssl_metrics}
\begin{tabular}{lccc}
\toprule
\textbf{Metric} & \textbf{Initial} & \textbf{Final} & \textbf{Improvement} \\
\midrule
Masked reconstruction MSE & -- & -- & -- \\
Ranking accuracy & -- & -- & -- \\
Cross-view correlation & -- & -- & -- \\
Spatial smoothness & -- & -- & -- \\
\bottomrule
\end{tabular}
\end{table}

\textbf{[RESULTS: Report convergence behavior, final loss values, and validation metrics. Describe whether overfitting occurred. Note any observations about learned representations prior to downstream tasks.]}


\subsection{Continuous Flow Matching with Optimal Transport}

The continuous flow matching transition model was evaluated on its ability to learn biologically coherent dynamics between adjacent disease stages. Training convergence was assessed via Sinkhorn divergence between predicted and ground-truth optimal transport couplings \cite{Cuturi2013-li}. Trajectory quality was measured by path straightness, comparing learned flows to linear interpolation baselines \cite{Lipman2022-ns}. Biological plausibility was validated by examining gene expression dynamics and pathway activation patterns along predicted trajectories.

\textbf{[FIGURE 3: CFM-OT transition model]}

\textit{Figure 3. Continuous flow matching results. (A) Learned vector field visualization in 2D UMAP projection. (B) Sinkhorn divergence over training epochs. (C) Trajectory straightness metrics comparing learned paths to linear interpolation. (D) Stage transition probability matrix derived from flow integration.}

\begin{table}[h]
\centering
\caption{CFM-OT transition metrics for each stage-to-stage transition.}
\label{tab:cfm_metrics}
\begin{tabular}{lccc}
\toprule
\textbf{Transition} & \textbf{Sinkhorn Div.} & \textbf{Path Straightness} & \textbf{Bio. Plausibility} \\
\midrule
Normal $\to$ AAH & -- & -- & -- \\
AAH $\to$ AIS & -- & -- & -- \\
AIS $\to$ MIA & -- & -- & -- \\
MIA $\to$ LUAD & -- & -- & -- \\
\bottomrule
\end{tabular}
\end{table}

\textbf{[FIGURE 4: Transition analysis]}

\textit{Figure 4. Biological analysis of learned transitions. (A) Gene expression dynamics along predicted trajectories for key marker genes. (B) Pathway activity scores (PROGENy) along trajectories. (C) Cell type composition shifts across transitions. (D) Predicted versus observed intermediate states.}

\textbf{[RESULTS: Report Sinkhorn divergence values, path straightness metrics, and biological validation of learned transitions. Describe gene expression dynamics and pathway activation patterns along predicted trajectories.]}


\subsection{Architecture Baseline Comparison}

To validate the architectural contributions of StageBridge, the full model was compared against four baselines representing progressively richer structural inductive biases: PoolingMLP (no structure), DeepSets (permutation invariance), SetTransformer (attention without spatial structure), and GraphSAGE (symmetric spatial message passing). All models were trained with identical hyperparameters and evaluated on held-out donors using five-fold cross-validation. Performance was assessed across reconstruction quality (MSE), representation quality (silhouette score, ARI), progression discrimination (AUROC), and batch integration (kBET).

\begin{table}[h]
\centering
\caption{Baseline comparison results. All metrics computed on held-out donors via 5-fold CV. $\downarrow$ indicates lower is better; $\uparrow$ indicates higher is better. Values shown as mean $\pm$ std.}
\label{tab:baselines}
\begin{tabular}{lccccc}
\toprule
\textbf{Model} & \textbf{MSE} $\downarrow$ & \textbf{Silhouette} $\uparrow$ & \textbf{ARI} $\uparrow$ & \textbf{AUROC} $\uparrow$ & \textbf{kBET} $\uparrow$ \\
\midrule
PoolingMLP & -- & -- & -- & -- & -- \\
DeepSets & -- & -- & -- & -- & -- \\
SetTransformer & -- & -- & -- & -- & -- \\
GraphSAGE & -- & -- & -- & -- & -- \\
\textbf{StageBridge} & -- & -- & -- & -- & -- \\
\bottomrule
\end{tabular}
\end{table}

\textbf{[FIGURE 5: Baseline comparison bar chart]}

\textit{Figure 5. Quantitative comparison of StageBridge against baselines. (A) Reconstruction MSE. (B) Silhouette score for stage clustering. (C) AUROC for stage classification. (D) kBET score for batch integration. Error bars show standard deviation across CV folds.}

\textbf{[RESULTS: Report metrics for each baseline. Discuss which architectural components explain performance differences. Identify which transitions from the baseline ladder produced the largest gains.]}


\subsection{Ablation Studies}

Systematic ablations were conducted to quantify the contribution of each architectural component to overall model performance. Reference ablations assessed the necessity of dual-reference geometry by comparing against single-reference variants. Spatial ablations removed the ring structure to isolate the contribution of explicit distance discretization. Receiver ablations replaced asymmetric attention with symmetric message passing. Loss ablations removed each auxiliary objective individually to quantify regularization benefits. All ablations were evaluated using the same cross-validation protocol as the baseline comparison.

\begin{table}[h]
\centering
\caption{Ablation study results. $\Delta$ columns show change relative to full model.}
\label{tab:ablations}
\begin{tabular}{lcccc}
\toprule
\textbf{Ablation} & \textbf{MSE} & $\Delta$\textbf{MSE} & \textbf{Silhouette} & $\Delta$\textbf{Sil.} \\
\midrule
Full model & -- & -- & -- & -- \\
\midrule
\multicolumn{5}{l}{\textit{Reference ablations}} \\
$-$ Dual reference (HLCA only) & -- & -- & -- & -- \\
$-$ Dual reference (LuCA only) & -- & -- & -- & -- \\
\midrule
\multicolumn{5}{l}{\textit{Architecture ablations}} \\
$-$ Spatial rings (flat attention) & -- & -- & -- & -- \\
$-$ Receiver centering (symmetric) & -- & -- & -- & -- \\
\midrule
\multicolumn{5}{l}{\textit{Loss ablations}} \\
$-$ Ranking loss & -- & -- & -- & -- \\
$-$ Consistency loss & -- & -- & -- & -- \\
$-$ Spatial loss & -- & -- & -- & -- \\
$-$ Group loss & -- & -- & -- & -- \\
\bottomrule
\end{tabular}
\end{table}

\textbf{[FIGURE 6: Ablation importance visualization]}

\textit{Figure 6. Contribution of architectural components. (A) Relative change in MSE for each ablation. (B) Relative change in silhouette score. Components ordered by importance.}

\textbf{[RESULTS: Report ablation metrics. Discuss which components contribute most. Interpret what this reveals about the importance of dual-reference geometry, spatial structure, and receiver-centering for progression modeling.]}


\subsubsection{Fusion Strategy Ablations}

To evaluate the contribution of dual-reference fusion design, we compared concatenation (default), learned weighted fusion, and confidence-weighted fusion strategies.

\begin{table}[h]
\centering
\caption{Fusion strategy ablation results. All metrics computed on held-out donors.}
\label{tab:fusion_ablations}
\begin{tabular}{lcccc}
\toprule
\textbf{Fusion Strategy} & \textbf{MSE} $\downarrow$ & \textbf{Silhouette} $\uparrow$ & \textbf{AUROC} $\uparrow$ & \textbf{ECE} $\downarrow$ \\
\midrule
Concatenation (default) & -- & -- & -- & -- \\
Learned weighted ($w=0.5$ init) & -- & -- & -- & -- \\
Confidence-weighted & -- & -- & -- & -- \\
\bottomrule
\end{tabular}
\end{table}

\textbf{[RESULTS: Report metrics for each fusion strategy. Discuss whether learned or confidence weighting improved over simple concatenation. Note the learned weight value at convergence---did it favor HLCA or LuCA?]}


\subsubsection{Loss Weight Ablations}

The SSL loss weight distribution was evaluated to test whether the 70\% emphasis on masked reconstruction was beneficial.

\begin{table}[h]
\centering
\caption{SSL loss weight ablation results.}
\label{tab:loss_weight_ablations}
\begin{tabular}{lcccc}
\toprule
\textbf{Weight Configuration} & \textbf{Recon MSE} $\downarrow$ & \textbf{Rank Acc} $\uparrow$ & \textbf{Silhouette} $\uparrow$ \\
\midrule
Default (70/10/10/5/5) & -- & -- & -- \\
Equal (20/20/20/20/20) & -- & -- & -- \\
Masked-only (100/0/0/0/0) & -- & -- & -- \\
\bottomrule
\end{tabular}
\end{table}

\textbf{[RESULTS: Report whether the 70\% masked reconstruction weight outperformed alternatives. Discuss the role of auxiliary losses in regularization---did removing them lead to overfitting or degraded representation quality?]}


\subsubsection{Confidence Calibration Ablations}

Temperature scaling was evaluated for its impact on confidence reliability.

\begin{table}[h]
\centering
\caption{Confidence calibration results. ECE measured on held-out validation set.}
\label{tab:calibration_ablations}
\begin{tabular}{lccc}
\toprule
\textbf{Calibration} & \textbf{ECE (HLCA)} $\downarrow$ & \textbf{ECE (LuCA)} $\downarrow$ & \textbf{Temperature $T$} \\
\midrule
None (raw percentile) & -- & -- & 1.0 (fixed) \\
Temperature scaling & -- & -- & -- (learned) \\
\bottomrule
\end{tabular}
\end{table}

\textbf{[RESULTS: Report ECE before and after temperature scaling. Discuss whether calibrated confidence improved downstream weighted fusion or other confidence-dependent components.]}


\subsection{Representation Quality}

The quality of learned cell representations was assessed through dimensionality reduction visualization and quantitative clustering metrics. UMAP projections were generated to examine whether representations organized by disease stage while preserving cell type identity. Batch integration was evaluated by coloring embeddings by donor to detect patient-specific clustering artifacts. Reference confidence scores were visualized to assess uncertainty distribution across the embedding space.

\textbf{[FIGURE 7: UMAP visualizations]}

\textit{Figure 7. UMAP projections of learned representations. (A) Colored by disease stage showing progression structure. (B) Colored by cell type showing biological identity preservation. (C) Colored by donor showing batch integration. (D) Colored by reference confidence showing uncertainty distribution.}

\textbf{[RESULTS: Describe stage separation, cell type preservation, and batch integration observed in learned embeddings. Quantify mixing using LISI scores.]}


\subsection{Attention Analysis}

The interpretability of learned attention patterns was examined to assess whether the model captured biologically meaningful sender-receiver relationships. Attention weights were aggregated across the test set and analyzed by spatial ring to characterize distance-dependent information flow. Cell type-specific attention patterns were computed by stratifying receivers into epithelial, stromal, and immune populations. Stage-specific patterns were examined to identify progression-dependent changes in niche interactions.

\textbf{[FIGURE 8: Attention weight analysis]}

\textit{Figure 8. Learned attention patterns. (A) Mean attention weight by spatial ring showing distance-dependent information flow. (B) Attention by sender cell type for epithelial receivers. (C) Attention by sender cell type for stromal receivers. (D) Stage-specific attention patterns showing progression-dependent changes.}

\textbf{[RESULTS: Report distance-dependent attention weighting. Describe cell type-specific and stage-specific attention patterns. Interpret biological meaning of learned attention---e.g., do epithelial cells attend more to fibroblasts at the invasion transition?]}


\subsection{Biological Validation}

Biological relevance of the learned representations was validated through multiple orthogonal analyses. Pathway activity scores were computed using PROGENy for both original and reconstructed expression profiles to assess preservation of functional signals. Marker gene expression was compared across cell types to verify that reconstructions maintained cell identity signatures. Gene set enrichment analysis was performed on genes with highest attention weights to characterize the biological processes captured by the model. Learned representations were compared against established EMT and cancer-associated fibroblast (CAF) gene signatures.

\textbf{[FIGURE 9: Biological validation]}

\textit{Figure 9. Biological validation. (A) Pathway activity correlation between original and reconstructed profiles (PROGENy). (B) Marker gene expression preservation by cell type. (C) Gene set enrichment of attention-weighted genes. (D) Correlation with established EMT and CAF signatures.}

\textbf{[RESULTS: Report pathway activity correlations (Pearson $r$). Quantify marker gene preservation by cell type. Describe enriched pathways from attention analysis---which biological processes does the model attend to?]}


\subsection{Progression Discrimination}

The ability of learned representations to discriminate disease progression stages was evaluated through classification and trajectory analysis. Pairwise binary classification was performed between adjacent stages (Normal vs AAH, AAH vs AIS, AIS vs MIA, MIA vs LUAD) using logistic regression on the learned embeddings. Five-way multiclass classification assessed overall stage discrimination. Pseudotime analysis using diffusion pseudotime was conducted to evaluate whether the embedding geometry recovered the expected histopathological ordering.

\textbf{[FIGURE 10: Progression analysis]}

\textit{Figure 10. Progression discrimination. (A) ROC curves for pairwise stage classification. (B) Five-way classification confusion matrix. (C) Diffusion pseudotime trajectory colored by stage. (D) Stage transition probabilities from Markov chain analysis.}

\textbf{[RESULTS: Report AUROC for each pairwise stage comparison (which transitions are hardest to discriminate?). Describe confusion matrix patterns. Report pseudotime trajectory analysis results---does embedding geometry recover the histopathological ordering?]}


\subsection{Computational Efficiency}

Computational requirements were measured for each stage of the StageBridge pipeline to characterize scalability and resource demands. Wall-clock runtime was recorded for reference mapping via scArches surgical fine-tuning, self-supervised pretraining, and CFM-OT transition model training. All experiments were conducted on NVIDIA H100 GPUs. Inference time was measured per cell to assess deployment feasibility for large-scale single-cell atlases.

\begin{table}[h]
\centering
\caption{Computational requirements for each pipeline component.}
\label{tab:compute}
\begin{tabular}{lccc}
\toprule
\textbf{Component} & \textbf{Time} & \textbf{Epochs} & \textbf{Hardware} \\
\midrule
LuCA scANVI retraining & $\sim$6 hrs & 329 + 34 & 4$\times$ H100 GPU \\
ENSG gene mapping & $<$5 min & -- & CPU \\
Reference mapping (HLCA) & $\sim$48 min & 60--61 & 1$\times$ H100 GPU \\
Reference mapping (LuCA) & TBD & TBD & 1$\times$ H100 GPU \\
SSL pretraining & -- & -- & 4$\times$ H100 GPU \\
CFM-OT training & -- & -- & 4$\times$ H100 GPU \\
Inference (per cell) & -- & -- & Single H100 \\
\bottomrule
\end{tabular}
\end{table}

LuCA scANVI retraining on 892K cells completed in approximately 6 hours on 4$\times$ H100 GPUs using PyTorch Lightning's \texttt{ddp\_spawn} distributed strategy. Early stopping triggered at epoch 329 for scVI (patience=20, monitoring validation ELBO) and epoch 34 for scANVI (patience=15), indicating efficient convergence without overfitting.

\textbf{[RESULTS: Report actual runtime for remaining pipeline components. Note early stopping behavior and convergence epochs. Discuss scalability to larger datasets.]}

